\clearpage
\VerdierPhysicalPageAnchor{0150}%
\begingroup
\small
\setlength{\parskip}{0pt}%
\setlength{\emergencystretch}{1.5em}%

\VerdierSemanticAnchor{ch2-2-3-5-w-continuation}
\noindent $B$ は $\mathcal B$ の対象であり、$s$ は
$S_{\mathcal D}(\mathcal A\cap\mathcal B)$ に属するので、$W$ は
$\mathcal B$ のある対象と同型である (1.2.4)。
\VerdierSemanticAnchor{ch2-2-3-5-i-implies-ii-conclusion}
したがって第一の主張により、$f$ は $\mathcal A\cap\mathcal B$ の対象を
経由して分解する。ゆえに $Q_3(f)$ は零であり、したがって $m=0$ である。
\VerdierSemanticAnchor{ch2-2-3-5-ii-implies-i}
次に ii) $\Rightarrow$ i) を示す。a) により、$B$ が $\mathcal B$ の対象、
$X$ が $\mathcal A$ の対象である $\mathcal D$ の任意の射
$g\colon B\longrightarrow X$ が $\mathcal A\cap\mathcal B$ の対象を経由して
分解することを示せばよい。
\VerdierSemanticAnchor{ch2-2-3-5-q3-g-zero}
仮定により
$Q_3(g)\in\operatorname{Hom}_{\mathcal D/(\mathcal A\cap\mathcal B)}(B,X)$
は零である。
\VerdierSemanticAnchor{ch2-2-3-5-t-morphism}
よって (2.2.4) により、$t\in S_{\mathcal D}(\mathcal A\cap\mathcal B)$
かつ $tg=0$ となる射 $t\colon X\longrightarrow X'$ が存在する。
\VerdierSemanticAnchor{ch2-2-3-5-distinguished-triangle}
$S_{\mathcal D}(\mathcal A\cap\mathcal B)$ の定義により、$B'$ が
$\mathcal A\cap\mathcal B$ の対象である形 $(B',X,X',i,t,w)$ の
$\mathcal D$ の完全三角が存在する。
\VerdierSemanticAnchor{ch2-2-3-5-proof-complete}
(1.2.1) により $tg=0$ から $g$ は $B'$ を経由して分解する。これで命題
(2.3.5) の証明が完了する。
\VerdierSemanticAnchor{ch2-2-3-5-left-localizing-definition}
\medskip
$\mathcal D^\circ$ の部分圏 $\mathcal A^\circ$ が右
$\mathcal B^\circ$-局所化であるとき、$\mathcal D$ の充満三角部分圏
$\mathcal A$ は\emph{左 $\mathcal B$-局所化}であるという。
\endgroup
